% Curs 9, pagina 6 din 8 — transcriere a manuscrisului din 1994.
% Fragment destinat adăugării la sfârşitul lui courses/course-09.tex,
% după pagina 5. Niciun proiect original nu a fost modificat.

\section*{Curs 9 - Pagina 6}

\noindent\underline{Exp.}

\noindent 1) $G=S_n$, \qquad $A=\{1,2,\ldots,n\}$.
\[
S_n\times A\longrightarrow A
\]
\noindent $(\sigma,i)\longmapsto \sigma i=\sigma(i)\in A$, \qquad
$\sigma:\{1,2,\ldots,n\}\longrightarrow\{1,2,\ldots,n\}$.

\bigskip
\noindent Se dem.\ că ac.\ aplic.\ e o acţ.

\bigskip
\noindent 2) $G$ gr.\ arbitrar, $A=G$.\footnote{Semnele celor două operaţii sunt greu de deosebit în fotografie; notaţia $*$/$\cdot$ redă distincţia explicită dintre operaţia externă şi cea internă.}
\[
\underset{\text{op. ext.}}{a*u}
\ \overset{\mathrm{def}}{=}\
\underset{\text{op. int.}}{a\cdot u}.
\]
% Notaţie editorială: simbolurile celor două operaţii sunt greu de
% deosebit în fotografie; este păstrată distincţia op. ext./op. int.

\bigskip
\noindent 3) $G$ grup, $A=G$.
\[
a*u\ \overset{\mathrm{def}}{=}\ aua^{-1}
\qquad\text{(acţ. prin conjugare)}.
\]
\noindent Se arată că e acţ.

\bigskip
\noindent 4) $G=\left\{
\underbrace{\begin{pmatrix}
\cos\alpha & -\sin\alpha\\
\sin\alpha & \cos\alpha
\end{pmatrix}}_{R_\alpha}
\ \middle|\ \alpha\in\mathbb{R}\right\}$.

\noindent $R_\alpha R_\beta=R_{\alpha\beta}$.\footnote{Indicele din dreapta pare a fi $\alpha\beta$; semnul ``$+$'' nu se distinge în fotografie. Relaţia uzuală pentru rotaţii este $R_\alpha R_\beta=R_{\alpha+\beta}$. Aici este păstrată lectura provizorie a manuscrisului.}

\bigskip
\noindent Dacă $\beta=-\alpha\ \Longrightarrow\
R_\alpha R_{-\alpha}=R_0=I_2=R_{-\alpha}R_\alpha$.

\bigskip
\noindent $G\leq GL_2(\mathbb{R})$.

\bigskip
\noindent Fie $A=\mathbb{R}^2$ şi not.\
$\begin{pmatrix}x\\y\end{pmatrix}\in\mathbb{R}^2$ (prin not.).

\[
\begin{pmatrix}
\cos\alpha & -\sin\alpha\\
\sin\alpha & \cos\alpha
\end{pmatrix}
\begin{pmatrix}x\\y\end{pmatrix}
\ \overset{\mathrm{def}}{=}\
\begin{pmatrix}
x\cos\alpha-y\sin\alpha\\
x\sin\alpha+y\cos\alpha
\end{pmatrix}.
\]

\bigskip
\noindent\underline{Obs.} Dacă $A$ e o $G$-mulţ.\ şi $H\leq G$
$\Longrightarrow A$ e şi $H$-mulţ.:
\[
H\times A\longrightarrow A.
\]

\noindent Dacă $\varphi:G\times A\longrightarrow A$, definim
$\varphi':H\times A\longrightarrow A$ prin
\[
\varphi'(a,u)=\varphi(a,u).
\]
\noindent $\varphi'$ satisf.\ (1) şi (2).

\bigskip
\noindent $\sigma\in S_n$ şi $H=\langle\sigma\rangle\leq S_n$.
\[
H\times\{1,2,\ldots,n\}\longrightarrow\{1,2,\ldots,n\}
\]
\noindent $H=\{\sigma^h\mid h\in\mathbb{Z}\}$.
\[
\sigma^h\cdot i=\sigma^h(i)\in\{1,\ldots,n\}.
\]
